% !TeX root = ../proyecto.tex

\begin{UseCase}{CU-05}{Gestión y Exportación de Seguimiento}{
	Describe el proceso mediante el cual el analista agrega una noticia investigada a la lista de "Seguimiento Activo" y posteriormente exporta dicho listado en un archivo CSV enriquecido para su uso en otras herramientas o reportes.
}
	\UCitem{Versión}{\color{Gray}
		1.0
	}\UCitem{Autor}{\color{Gray}
		Tesista de Ingeniería
	}\UCitem{Supervisa}{\color{Gray}
		Director de Trabajo Terminal
	}\UCitem{Actor}{
		Analista de Datos / Trabajador Social
	}\UCitem{Propósito}{\begin{Titemize}
		\Titem Facilitar el monitoreo continuo de los casos más críticos identificados por el sistema y proporcionar interoperabilidad mediante la exportación de datos tabulares.
	\end{Titemize}
	}\UCitem{Entradas}{\begin{Titemize}
		\Titem ID de la noticia (para agregar a seguimiento).
		\Titem Petición de exportación (clic en botón).
	\end{Titemize}
	}\UCitem{Origen}{\begin{Titemize}
		\Titem Interfaz de resultados de investigación y sección de "Seguimiento" en el Dashboard.
	\end{Titemize}
	}\UCitem{Salidas}{\begin{Titemize}
		\Titem Confirmación visual en el Dashboard.
		\Titem Archivo descargable en formato \texttt{.csv} con los casos enriquecidos.
	\end{Titemize}
	}\UCitem{Destino}{\begin{Titemize}
		\Titem Navegador del usuario (descarga de archivo).
	\end{Titemize}
	}\UCitem{Precondiciones}{\begin{Titemize}
		\Titem La noticia debe tener una investigación estructurada (\texttt{DeepInvestigator}) completada exitosamente.
	\end{Titemize}
	}\UCitem{Postcondiciones}{\begin{Titemize}
		\Titem La noticia cambia su estado lógico en la base de datos a "En Seguimiento".
	\end{Titemize}
	}\UCitem{Errores}{\begin{Titemize}
		\Titem {\bf \hypertarget{CU-05.E1}{E1}}: Intento de exportación cuando la lista de seguimiento está vacía \Trayref{CU-05}{A}.
	\end{Titemize}
	}\UCitem{Tipo}{
		Caso de uso primario
	}\UCitem{Observaciones}{
		La exportación integra de forma transparente los campos relacionales de la noticia original junto con los campos JSON extraídos por el módulo de investigación (edades, ubicación, contacto DIF).
	}
\end{UseCase}

%--------------------------------------
\begin{UCtrayectoria}
	\UCpaso[\UCActor] Analiza los resultados generados por el módulo \texttt{DeepInvestigator} para un caso específico.
	\UCpaso[\UCActor] Determina que el caso requiere monitoreo a largo plazo y hace clic en el botón \IUbutton{Agregar a Seguimiento}.
	\UCpaso[\UCsist] Actualiza la bandera de seguimiento (\texttt{is\_tracked = True}) para el registro en la base de datos.
	\UCpaso[\UCsist] Confirma la acción mediante un mensaje en la interfaz.
	\UCpaso[\UCActor] Navega hacia la vista de "Casos en Seguimiento" del Dashboard.
	\UCpaso[\UCActor] Hace clic en el botón \IUbutton{Exportar a CSV}.
	\UCpaso[\UCsist] Consulta la base de datos extrayendo todos los casos marcados en seguimiento \ErrorRef{CU-05}{E1}{Lista Vacía}.
	\UCpaso[\UCsist] Serializa los datos, desempacando el JSON de la ficha investigativa en columnas planas.
	\UCpaso[\UCsist] Transmite el archivo \texttt{seguimiento\_casos.csv} como una descarga HTTP al navegador.
	\UCpaso[\UCActor] Recibe y guarda el archivo en su computadora local.
\end{UCtrayectoria}

%--------------------------------------
\begin{UCtrayectoriaA}{CU-05}{A}{El analista solicita la exportación pero no hay casos en seguimiento}
	\UCpaso[\UCsist] Detecta que la consulta a la base de datos retorna cero registros.
	\UCpaso[\UCsist] Aborta la generación del archivo CSV.
	\UCpaso[\UCsist] Muestra una notificación visual (\textit{Toast} o \textit{Alert}) indicando "No hay casos en seguimiento para exportar".
	\UCpaso[] El Caso de Uso termina.
\end{UCtrayectoriaA}
